% ======================================================================
% Ontology of Continua -- Core 1.3
% Cross-K module: crossk_k7_k8.tex
% Cross-level relations between K7 and K8
% ======================================================================

\subsubsection{Межуровневая структура K7–K8}
\label{sec:crossk-k7-k8}

Переход $K_7 \rightarrow K_8$ обозначает возникновение
цивилизационно-технологического континуума из социального континуума.
Если $K_7$ даёт нормы, роли, коммуникационные потоки и институциональные
циклы, то $K_8$ вводит:
\begin{itemize}
  \item технологические и инфраструктурные оси,
  \item крупномасштабные символические и информационные системы,
  \item сети энергии и логистики,
  \item экономические механизмы воспроизводства,
  \item циклы ремонта, обслуживания и высокоуровневой кооперации,
  \item глобальные условия непрерывности для систем популяционного масштаба.
\end{itemize}

$K_8$ — первый континуум, способный поддерживать устойчивые структуры,
превышающие когнитивные и социальные горизонты отдельного агента.

% ----------------------------------------------------------------------
\subsubsection{1. Задействованные уровни и их роли}

\paragraph{K7 (социальный континуум).}
$K_7$ предоставляет:
\begin{itemize}
  \item социальные нормы и роли ($A_{\mathrm{soc}}$),
  \item коммуникационные потоки $J_{\mathrm{comm}}$,
  \item институциональные циклы $C_{\mathrm{inst}}$,
  \item коллективные потенциалы $P_{\mathrm{soc}}$ (доверие, легитимность),
  \item пороги сплоченности и принуждения $\Theta_{\mathrm{soc}}$.
\end{itemize}

Однако в $K_7$ отсутствуют:
\begin{itemize}
  \item инфраструктуры большого масштаба,
  \item технологические циклы производства,
  \item пороги плотности энергии,
  \item материально-информационные логистики,
  \item устойчивые механизмы многопоколенного воспроизводства.
\end{itemize}

\paragraph{K8 (цивилизационно-технологический континуум).}
$K_8$ вводит:
\begin{itemize}
  \item оси для инфраструктуры, технологий и символических систем
        ($A_{\mathrm{tech}}$),
  \item потенциалы энергоёмкости, ремонтопригодности, пропускной
        способности логистики,
  \item потоки $J^{(8)} = (J_{\mathrm{energy}}, J_{\mathrm{material}},
        J_{\mathrm{logistics}}, J_{\mathrm{info}}, J_{\mathrm{population}})$,
  \item инфраструктурные циклы $C_{\mathrm{sys}}$
        (производство → транспорт → потребление → обслуживание),
  \item пороги $\Theta_8 = (\Theta_E, \Theta_L, \Theta_I, \Theta_M, \Theta_R)$,
  \item меру $k_8$, описывающую цивилизационную непрерывность.
\end{itemize}

Так $K_8$ становится первым континуумом с крупномасштабной энергетической
и информационной автономностью.

% ----------------------------------------------------------------------
\subsubsection{2. Совместные и наследуемые оси и пороги}

\paragraph{Общие символические структуры.}
Символические структуры и коммуникационные системы $K_7$ служат основой для:
\[
A_{\mathrm{symbol}}(K_7) \rightarrow A_{\mathrm{tech}}(K_8),
\]
что делает возможной техническую кодировку (письмо, математика, инженерия,
стандартизация).

\paragraph{Новые оси в \texorpdfstring{$K_8$}{K_8}.}
Сюда входят:
\begin{itemize}
  \item $A_{\mathrm{infra}}$ — дороги, энергосети, водные системы,
  \item $A_{\mathrm{energy}}$ — энергоёмкость и системы преобразования энергии,
  \item $A_{\mathrm{econ}}$ — экономические механизмы воспроизводства,
  \item $A_{\mathrm{info}}$ — хранение и передача информации,
  \item $A_{\mathrm{repair}}$ — структуры обслуживания и регенерации.
\end{itemize}

\paragraph{Пороги.}
Наследованные социальные пороги сплоченности остаются необходимыми, но
недостаточными.
Появляются новые пороги:
\begin{itemize}
  \item $\Theta_E$ — минимальная энергоёмкость для воспроизводства на уровне
        системы,
  \item $\Theta_L$ — пороги пропускной способности логистики,
  \item $\Theta_I$ — пороги информационной когерентности,
  \item $\Theta_M$ — пороги ремонтного потенциала,
  \item $\Theta_R$ — пороги демографической устойчивости.
\end{itemize}

При нарушении любого из них $\Omega(K_8)$ схлопывается.

% ----------------------------------------------------------------------
\subsubsection{3. Межуровневые потоки, циклы и напряжения}

\paragraph{Потоки.}
Потоки в $K_8$ обобщают и масштабируют социальные потоки:
\[
J^{(8)} = (J_{\mathrm{energy}}, J_{\mathrm{material}},
          J_{\mathrm{logistics}}, J_{\mathrm{info}},
          J_{\mathrm{population}}).
\]

\paragraph{Циклы.}
Цивилизационная устойчивость требует инфраструктурного цикла $C_{\mathrm{sys}}$:
\[
\text{production} \rightarrow \text{transport} \rightarrow
\text{use} \rightarrow \text{maintenance} \rightarrow \text{production}.
\]
Условие устойчивости:
\[
\oint dR \approx 0.
\]

\paragraph{Напряжение.}
Межуровневое напряжение:
\[
T_{7,8} =
f(\Theta_E - P_{\mathrm{energy}},\
  \Theta_L - J_{\mathrm{logistics}},\
  \Theta_I - P_{\mathrm{info}},\
  \Theta_M - P_{\mathrm{repair}},\
  \Theta_R - P_{\mathrm{population}}).
\]

Избыточное напряжение разрушается инфраструктурную непрерывность.

% ----------------------------------------------------------------------
\subsubsection{4. Условия рождения и исчезновения между уровнями}

\paragraph{Рождение \texorpdfstring{$K_8$}{K_8}.}
$K_8$ появляется, когда:
\[
T_7 > \Theta_{7,\mathrm{dim}}
\quad\text{и}\quad
A_{\mathrm{tech}}, A_{\mathrm{infra}} \in M_8 \setminus A(K_7).
\]

Домен расширяется:
\[
\Omega(K_8) = \Omega(K_7) \cup \Delta\Omega_{\mathrm{civ}}.
\]

\paragraph{Распространение гибели.}
Если социальные пороги нарушаются ($\Theta_{\mathrm{soc}}$), цивилизационные
системы теряют когерентность и схлопываются в $K_7$.

Если нарушается любой из пяти порогов $K_8$:
\[
\Theta_E,\Theta_L,\Theta_I,\Theta_M,\Theta_R,
\]
то $\Omega(K_8)=\varnothing$.

$K_7$ может оставаться устойчивой, пока его собственные пороги не нарушены.

% ----------------------------------------------------------------------
\subsubsection{5. Концептуальные примеры}

\paragraph{Рост энергоёмкости.}
Когда энергоёмкость превышает $\Theta_E$, технологические и экономические
структуры становятся самоподдерживающимися.

\paragraph{Информационная кодировка.}
Социальные символы переходят в технические системы (писмо → математика →
наука), делая возможным появление новых осей в $A_{\mathrm{tech}}$.

\paragraph{Сбой инфраструктуры.}
Если емкость логистики снижается ниже $\Theta_L$ или темпы ремонта ниже
$\Theta_M$, цивилизационные циклы разрываются, и система коллапсирует:
\[
C_{\mathrm{sys}} \nrightarrow \text{замкнутый цикл}.
\]

Следовательно, переход K7→K8 в сжатой форме — это возникновение
технологически-цивилизационной структуры из социальных норм, коммуникационных
потоков и институциональной устойчивости.
